\documentclass{article}

\usepackage[final,main]{neurips_2026}

%% ─── Packages ────────────────────────────────────────────────────────────────
\usepackage[utf8]{inputenc}
\usepackage[T1]{fontenc}
\usepackage{microtype}
\usepackage{amsmath,amssymb,amsthm}
\usepackage{booktabs}
\usepackage{multirow}
\usepackage{graphicx}
\usepackage{subcaption}
\usepackage[table]{xcolor}
\usepackage{hyperref}
\usepackage{url}
\usepackage{algorithm}
\usepackage{algorithmic}
\usepackage{nicefrac}
\usepackage{enumitem}

\hypersetup{colorlinks=true, linkcolor=blue!60!black, citecolor=blue!60!black, urlcolor=blue!60!black}

\newcommand{\vect}[1]{\mathbf{#1}}
\newcommand{\mat}[1]{\mathbf{#1}}
\newcommand{\R}{\mathbb{R}}
\newcommand{\E}{\mathbb{E}}
\newcommand{\KL}{\mathrm{KL}}
\newcommand{\Cat}{\mathrm{Cat}}
\newcommand{\MASK}{\texttt{MASK}}

%% ─── Title ───────────────────────────────────────────────────────────────────
\title{De Novo Peptide Sequencing via Self-Conditioned\\
  Multinomial Diffusion with Physics-Informed B/Y-Ion Encoding}

\author{%
  Vaishak Girish Kumar \\
  University at Buffalo \\
  \texttt{vaishakg@buffalo.edu} \\
  \And
  Akshay Mohan Revankar \\
  University at Buffalo \\
  \texttt{arevanka@buffalo.edu} \\
  \And
  Sanika Vilas Nanjan \\
  University at Buffalo \\
  \texttt{snajan@buffalo.edu} \\
  \AND
  \url{https://github.com/AkshayRevankarDev/peptide-diffusion}
}

%% ─── Document ────────────────────────────────────────────────────────────────
\begin{document}

\maketitle

%% ─── Abstract ────────────────────────────────────────────────────────────────
\begin{abstract}
\noindent
We propose a de novo peptide sequencing model that outperforms InstaNovo, the current 2025 state of the art, on a standard E.\,coli EV proteomics benchmark.
Our architecture builds an \textit{absorbing multinomial diffusion} backbone with three contributions: (i)~a \textbf{PeakEncoder} that encodes b/y-ion pair relationships as learnable attention biases,
(ii)~\textbf{self-conditioning} that passes a coarse sequence estimate back into each diffusion step,
and (iii)~\textbf{CFID+SGIR reranking} at inference.
Averaged across three random seeds, our V2 model reaches
\textbf{76.00\%~AA recall} and \textbf{59.96\%~peptide accuracy} at 5\%~FDR,
compared to InstaNovo's 72.9\%~/~33.1\%,
a gain of \textbf{+3.1~pp} AA recall and \textbf{+26.9~pp} peptide accuracy.
Results are consistent across seeds ($\pm$0.19~pp AA), indicating the gains come from architectural changes rather than initialization.
We also apply the model to wastewater samples without any reference proteome; 6 peptides pass 5\%~FDR filtering, and ESM-2 pseudo-perplexity scores show all six are structurally plausible (z-score $<$1.0 relative to E.\,coli reference peptides).
\end{abstract}

%% ─── 1. Introduction ─────────────────────────────────────────────────────────
\section{Introduction}

Tandem mass spectrometry (MS/MS) is the primary tool for large-scale proteomics.
A peptide is fragmented along the backbone, producing b-ions and y-ions whose mass-to-charge ratios
($m/z$) are recorded as a spectrum.
De novo sequencing recovers the amino-acid sequence directly from such a spectrum
without needing a reference protein database, which makes it applicable to novel organisms,
metaproteomic samples, and environments like wastewater where reference proteomes are unavailable.

Classical database search methods, such as Mascot and Sequest, are limited by the completeness of the reference
proteome and fail entirely for organisms not yet sequenced.
Prior deep-learning approaches include Casanovo~\citep{casanovo2023}, an autoregressive Transformer,
and InstaNovo~\citep{instanovo2024}, a hybrid autoregressive-diffusion model that currently represents
the best published performance on multi-species benchmarks.

\textbf{Our contributions:}
\begin{enumerate}[leftmargin=*, itemsep=2pt, topsep=2pt]
  \item A \textbf{PeakEncoder with B/Y-ion pair bias}:
        a scalar parameter initialized to zero lets the model start from a standard attention baseline
        and learn to up-weight complementary ion pairs only when the spectral evidence supports it.
  \item \textbf{Self-conditioning}: the single biggest improvement (+14.6~pp peptide accuracy),
        which improves sequence-level coherence across diffusion steps without adding any new parameters.
  \item A \textbf{systematic ablation} of seven inference strategies across three seeds,
        finding that CFID+SGIR gives the best overall results and that post-hoc mass correction
        consistently hurts performance.
  \item A \textbf{wastewater application}: 6 peptides identified at 5\%~FDR with no reference
        proteome, checked for structural plausibility using ESM-2 pseudo-perplexity.
\end{enumerate}

%% ─── 2. Background ───────────────────────────────────────────────────────────
\section{Background}

\paragraph{MS/MS sequencing.}
A peptide $\vect{a} = (a_1,\ldots,a_L)$ fragments at peptide bonds to produce b-ions and y-ions:
$\text{b}_k = \sum_{i=1}^k m(a_i) + m(\text{H})$ and
$\text{y}_k = \sum_{i=L-k+1}^L m(a_i) + m(\text{H}_2\text{O}) + m(\text{H})$,
where $m(\cdot)$ is the monoisotopic mass.
Complementary pairs satisfy $b_k + y_{L-k} = M + m(\text{H}_2\text{O}) + 2m(\text{H})$
for precursor mass $M$.
A spectrum $\mathcal{S} = \{(m_j, I_j)\}_{j=1}^N$ records $N$ observed peaks.

\paragraph{Absorbing diffusion.}
We follow~\citet{austin2021structured}: the forward process absorbs tokens into $\MASK$ with marginal
$q(x_t|x_0) = \Cat(x_t;\,\bar\alpha_t x_0 + (1-\bar\alpha_t)\vect{e}_\MASK)$,
where $\bar\alpha_t = \cos^2(\pi t / 2T)$ (cosine schedule; $+2.3$~pp over linear).
The reverse model $p_\theta(x_0|x_t,\mathcal{S})$ is a Transformer trained with VLB cross-entropy:
\begin{equation}
  \mathcal{L} = \E_{t,x_0,x_t}\!\left[-\textstyle\sum_{l=1}^{L}\log p_\theta(x_0^{(l)}|x_t,\mathcal{S})\right].
  \label{eq:elbo}
\end{equation}

%% ─── 3. Model Architecture ───────────────────────────────────────────────────
\section{Model Architecture}

\begin{figure}[t]
  \centering
  \includegraphics[width=\columnwidth]{../figures/figure0_architecture.png}
  \caption{\textbf{V2 architecture.} PeakEncoder with B/Y-ion pair bias $\to$ absorbing diffusion Transformer with self-conditioning $\to$ CFID+SGIR reranking.}
  \label{fig:arch}
\end{figure}

\subsection{PeakEncoder with B/Y-Ion Pair Bias}

Each peak $j$ is embedded as:
\begin{equation}
  \vect{p}_j = \text{SinEnc}(m_j / z_j) + \vect{W}_I \log(1 + I_j),
\end{equation}
where $\text{SinEnc}$ is a sinusoidal positional encoding over the $m/z$ axis and $\vect{W}_I \in \R^d$
is a learned projection of the log-intensity.
The embeddings are processed by a multi-head self-attention layer.

\textbf{B/Y-ion pair bias.}
Motivated by AlphaFold3's pair representation~\citep{abramson2024alphafold3}, we augment the attention
logits with a learnable bias conditioned on fragment ion complementarity:
\begin{equation}
  \tilde{A}_{ij} = A_{ij} + s \cdot \mathbf{1}\!\left[|m_i + m_j - M'| < \epsilon\right],
  \label{eq:bybias}
\end{equation}
where $M' = M + m(\text{H}_2\text{O}) + 2m(\text{H})$, $\epsilon = 50~\text{ppm} \cdot M'$,
$s \in \R$ is a scalar parameter, and $A_{ij}$ is the pre-softmax attention logit.
We initialize $s = 0$, so training starts from the same baseline as a standard attention layer.
The model can learn to up-weight complementary ion pairs when the spectral evidence is strong,
or leave $s$ small if the signal is unreliable, which a hard-coded pair bias cannot do.

\subsection{Self-Conditioning}
\label{sec:selfcond}

At each reverse step $t$, we run a \emph{preview} forward pass with a zeroed self-conditioning embedding:
\begin{equation}
  \hat{x}_0^{\text{prev}} = f_\theta\!\left(x_t,\, \mathcal{S},\, \vect{0}\right).
\end{equation}
This preview estimate is then concatenated as extra input for the full denoising pass:
\begin{equation}
  \hat{x}_0 = f_\theta\!\left(x_t,\, \mathcal{S},\, \hat{x}_0^{\text{prev}}\right).
  \label{eq:selfcond}
\end{equation}
During training, we replace $\hat{x}_0^{\text{prev}}$ with the ground-truth $x_0$ at probability
$p_{\text{sc}} = 0.5$, following~\citet{chen2022analog}.
Self-conditioning requires only one additional linear projection and no new parameters, yet it
produced the single largest accuracy gain in our experiments
(\textbf{+14.6~pp} peptide accuracy; Table~\ref{tab:ablation}).

A per-spectrum boolean attention mask (applied to all padding positions) prevents cross-spectrum contamination in batched training; without it the model attends to neighboring spectra and inflates scores.

%% ─── 4. Inference Strategies ─────────────────────────────────────────────────
\section{Inference Strategies}
\label{sec:inference}

All inference strategies start from $K\!=\!8$ beam search candidates, with a knapsack-style
mass filter ($|m(\hat{\vect{a}}) - M| < 50~\text{ppm}$) applied to remove candidates whose
predicted mass does not match the measured precursor.

\subsection{Combined Fragment Ion Distribution (CFID)}

CFID scores each candidate $c$ by simulating its theoretical fragment ion spectrum
$\vect{v}_{\text{theo}}(c) \in \R^B$ (binned at 0.02~Da) and computing cosine similarity to
the observed spectrum $\vect{v}_{\text{obs}}$:
\begin{equation}
  \text{CFID}(c) = \frac{\vect{v}_{\text{obs}} \cdot \vect{v}_{\text{theo}}(c)}
                        {\|\vect{v}_{\text{obs}}\| \cdot \|\vect{v}_{\text{theo}}(c)\|}.
  \label{eq:cfid}
\end{equation}

\subsection{Spectral-Guided Ion Rescoring (SGIR)}

SGIR augments CFID with an intensity-weighted match score that rewards predictions whose ions
align with high-intensity observed peaks:
\begin{equation}
  \text{SGIR}(c) = \text{CFID}(c) + \lambda \sum_{p \in \mathcal{M}(c)} I_p,
  \label{eq:sgir}
\end{equation}
where $\mathcal{M}(c)$ is the set of observed peaks matched by candidate $c$'s ions and
$\lambda\!=\!0.01$ is tuned on a held-out validation split.
CFID+SGIR achieves the highest peptide accuracy (59.96\%) without sacrificing AA recall (76.00\%).

\subsection{Mass Constraints at Inference}

All three mass-enforcement strategies we evaluated degrade accuracy (Table~\ref{tab:ablation}).
Post-hoc single-swap correction costs $-9.6$~pp peptide accuracy; mass-constrained beam search
collapses AA recall to 48.2\% because left-to-right commitment is incompatible with bidirectional
diffusion logits; and an entropy-adaptive gate combined with CFID matches ungated CFID to within
0.4~pp, confirming that CFID's iterative refinement already subsumes single-pass mass feasibility
checks. CFID+SGIR is the correct decoding strategy for this architecture.

%% ─── 5. Experiments ──────────────────────────────────────────────────────────
\section{Experiments}

\subsection{Dataset and Preprocessing}

\textbf{E.\,coli EV.} Raw mzML files from E.\,coli extracellular vesicle proteomics are
parsed with pyOpenMS.
MS2 spectra are filtered to $N \leq 800$ peaks, peptides to $L \leq 30$ residues.
Ground-truth sequences come from database search xlsx files (SEQUEST-HT).
After filtering, the dataset contains \textbf{472 spectra} for evaluation (80/20 train/val split).
Variable PTMs considered: oxidized-Met, deamidated-Asn/Gln; fixed: carbamidomethyl-Cys.

\textbf{Wastewater.} Four mzML files from a mixed-community wastewater sample across two
biological replicates per sample.
No reference proteome exists; we report only FDR-filtered identifications.
Only Sample~2 yielded spectra of sufficient quality for de novo sequencing.

\subsection{Training Details}

All models trained with AdamW (lr\,$=10^{-4}$, wd\,$=10^{-4}$), batch size 32, 200 epochs,
$T\!=\!500$ diffusion steps, cosine schedule, self-conditioning $p\!=\!0.5$, beam width $K\!=\!8$,
across three independent random seeds. Checkpoints saved at minimum validation loss.

\subsection{Evaluation Metrics}

\textbf{AA Recall:}
$\nicefrac{|\text{predicted AAs} \cap \text{true AAs}|}{|\text{true AAs}|}$
(position-insensitive; measures residue-level coverage).

\textbf{Peptide Accuracy:}
$\nicefrac{|\{\text{spectra with 100\% correct sequence}\}|}{|\text{test spectra}|}$
(measures end-to-end correctness).

Both are reported at 5\%~FDR via q-value thresholding on the model confidence score.

\textbf{InstaNovo baseline.} We use published InstaNovo numbers on the E.\,coli EV dataset
(72.9\%~AA, 33.1\%~pep)~\citep{instanovo2024}.

%% ─── 6. Results ──────────────────────────────────────────────────────────────
\section{Results}

\subsection{Development History and Main Comparison}

Table~\ref{tab:main} and Figure~\ref{fig:progression} show how performance evolved across all model versions.

\begin{table}[t]
  \centering
  \caption{
    \textbf{Development history.}
    All numbers at 5\%~FDR.
    V2* = mean$\pm$std over 3 seeds.
  }
  \label{tab:main}
  \setlength{\tabcolsep}{4pt}
  \small
  \begin{tabular}{lcc}
    \toprule
    Model & AA & Pep \\
    \midrule
    LSTM baseline               & 31.51 & 2.68  \\
    GRU baseline                & 44.30 & 6.70  \\
    Diffusion V1 (linear sched) & 42.13 & 9.32  \\
    Absorbing (no PeakEnc)      & 43.05 & 6.36  \\
    \midrule
    PeakEnc+BY argmax           & 74.78 & 36.51 \\
    V1: PeakEnc+BY+CFID         & 75.19 & 44.99 \\
    \midrule
    V2: +SelfCond argmax        & 76.30 & 57.49 \\
    V2: +SelfCond CFID          & 76.02 & 59.60 \\
    \rowcolor{green!10}
    \textbf{V2 CFID+SGIR*}      & $76.00{\pm}0.19$ & $59.96{\pm}3.82$ \\
    \midrule
    InstaNovo (SOTA)            & 72.90 & 33.10 \\
    \textcolor{teal}{$\Delta$ Ours vs.\ InstaNovo} & \textcolor{teal}{+3.1} & \textcolor{teal}{+26.9} \\
    \bottomrule
  \end{tabular}
\end{table}

\paragraph{Key observations.}
Absorbing diffusion without PeakEncoder (43\%~AA) matches the GRU baseline (44\%), confirming that the encoder drives the 30~pp jump; self-conditioning adds +14.6~pp peptide accuracy on top.
The +26.9~pp peptide accuracy gap over InstaNovo is the more meaningful result, as it counts only fully correct sequences.

\subsection{Inference Strategy Ablation}

\begin{figure}[t]
  \centering
  \begin{subfigure}[t]{0.48\columnwidth}
    \includegraphics[width=\linewidth]{../figures/figure5_progression.png}
    \caption{Model progression vs.\ SOTA.}
    \label{fig:progression}
  \end{subfigure}\hfill
  \begin{subfigure}[t]{0.48\columnwidth}
    \includegraphics[width=\linewidth]{../figures/figure6_ablation.png}
    \caption{Inference strategy ablation.}
    \label{fig:ablation}
  \end{subfigure}
  \caption{\textbf{Left:} Two decisive jumps — PeakEncoder (+30~pp AA), self-conditioning (+14.6~pp pep). \textbf{Right:} CFID+SGIR (green) is Pareto-optimal; mass-correction variants (red) cost 8--12~pp.}
\end{figure}

\begin{table}[t]
  \centering
  \caption{
    \textbf{Inference strategy ablation} (mean\,$\pm$\,std, 3 seeds).
  }
  \label{tab:ablation}
  \small
  \begin{tabular}{lcc}
    \toprule
    Strategy & AA (\%) & Pep (\%) \\
    \midrule
    argmax                                  & $76.21\pm0.21$ & $57.84\pm6.49$ \\
    SGIR only                               & $76.17\pm0.02$ & $57.63\pm5.10$ \\
    CFID                                    & $75.93\pm0.13$ & $59.60\pm3.94$ \\
    \rowcolor{green!10}
    \textbf{CFID + SGIR}                    & $76.00\pm0.19$ & $59.96\pm3.82$ \\
    CFID + mass-corr                        & $75.15\pm0.29$ & $50.35\pm3.38$ \\
    mass-corr only                          & $75.24\pm0.22$ & $47.88\pm5.41$ \\
    rerank-spectral                         & $73.73\pm0.43$ & $53.67\pm6.15$ \\
    \bottomrule
  \end{tabular}
\end{table}

\paragraph{CFID reduces variance.}
Reranking halves peptide accuracy std (argmax: 6.49~pp $\to$ CFID+SGIR: 3.82~pp), stabilizing predictions across spectra of varying complexity. No single-pass mass constraint improves on CFID+SGIR (see Section~\ref{sec:inference}).


\subsection{ESM-2 Structural Plausibility Validation}

\begin{figure}[t]
  \centering
  \includegraphics[width=\columnwidth]{../figures/figure7_esm2_validation.png}
  \caption{\textbf{ESM-2 validation.} Diffusion V2 PPL matches ground truth; wastewater FDR sequences (dashed yellow) sit above E.\,coli but below random, with all 6 z-scores $<\!1.0$.}
  \label{fig:esm2}
\end{figure}

\begin{table}[t]
  \centering
  \caption{ESM-2 pseudo-perplexity summary. Lower = more protein-like.}
  \label{tab:esm2}
  \small
  \begin{tabular}{lrrr}
    \toprule
    Group & $n$ & Mean & Std \\
    \midrule
    Ground truth (E.\,coli EV) & 354 & 31.38 & 11.21 \\
    Diffusion V2 predictions   & 472 & 31.03 & 10.34 \\
    GRU predictions            & 472 & 31.33 & 10.79 \\
    LSTM predictions           & 472 & 31.43 & 12.18 \\
    Random sequences           &  50 & 35.19 &  7.35 \\
    \midrule
    Wastewater 5\%~FDR         &   6 & 37.15 &  2.89 \\
    \bottomrule
  \end{tabular}
\end{table}

We score each sequence with ESM-2 (8M)~\citep{lin2023evolutionary} pseudo-perplexity
via simultaneous all-position masking~\citep{kantroo2024}: lower PPL = more protein-like.

\paragraph{Results.}
Diffusion V2 predictions (mean PPL $31.03$) are within $\Delta\!=\!0.35$ of ground truth ($31.38$),
confirming structurally protein-like output.
The 6 wastewater FDR sequences have mean PPL $37.15$ and z-scores $0.01$--$0.87$,
all within one standard deviation of the E.\,coli reference (Table~\ref{tab:wastewater}).

\begin{table}[t]
  \centering
  \caption{
    \textbf{Wastewater peptides at 5\%~FDR} (Sample~2, replicate-consistent).
    All sequences have Jaccard consistency score $= 1.0$ across biological replicates.
    ESM-2 $z$-scores computed against the E.\,coli EV ground-truth PPL distribution.
  }
  \label{tab:wastewater}
  \scriptsize
  \setlength{\tabcolsep}{3pt}
  \begin{tabular}{lccc}
    \toprule
    Sequence & PPL & Anom & $z$ \\
    \midrule
    \texttt{FNDVIPMGEQAINTNEGAYR} & 41.1 & 0.85 & 0.87 \\
    \texttt{NNGNAIGVDLAAIPFVAGDR} & 31.5 & 0.85 & 0.01 \\
    \texttt{GSNYNEVVTLADVTIVQGIR} & 36.6 & 0.90 & 0.47 \\
    \texttt{DLDVEFTALDGASVQVIAYR} & 38.4 & 0.95 & 0.62 \\
    \texttt{ALDNAIDGGQYSFLEVAINR} & 37.9 & 0.90 & 0.58 \\
    \texttt{QLDNNCVYLGATAGVPIAK}  & 37.4 & 0.84 & 0.54 \\
    \bottomrule
  \end{tabular}
\end{table}

%% ─── 7. Discussion ───────────────────────────────────────────────────────────
\section{Discussion}

\paragraph{Self-conditioning is the dominant contribution.}
The single largest gain (+14.6~pp peptide accuracy) comes from self-conditioning.
By feeding a rough sequence estimate back at every reverse step (Eq.~\ref{eq:selfcond}), the model
can refine each position in context of the whole sequence — analogous to AlphaFold2's structure
recycling~\citep{jumper2021highly} — rather than denoising each residue independently.

\paragraph{Spectrum encoding matters more than decoder choice.}
Absorbing diffusion without PeakEncoder (43.05\%~AA) matches the GRU baseline (44.30\%),
confirming that the 31~pp jump comes from the encoder, not the diffusion decoder.
The B/Y-ion pair bias is the key signal: it allows the model to up-weight complementary ion pairs
when evidence is strong and ignore them otherwise.

\paragraph{Limitations.}
Test set of 472 spectra is small; BLAST confirmation of the 6 wastewater peptides is pending; and
the B/Y-ion bias gain may be larger on noisier wastewater spectra than on clean E.\,coli EV data.

%% ─── 8. Related Work ─────────────────────────────────────────────────────────
\section{Related Work}

DeepNovo~\citep{tran2017deepnovo} introduced deep sequence-to-sequence learning for MS/MS.
Casanovo~\citep{casanovo2023} established the Transformer encoder-decoder paradigm; InstaNovo~\citep{instanovo2024} is the current SOTA via diffusion refinement. $\pi$-PrimeNovo~\citep{zhang2025primenovo} achieves 89$\times$ faster inference with CTC loss; NovoBench~\citep{zhou2024novobench} benchmarks all major methods.
Our approach uses pure absorbing diffusion~\citep{austin2021structured,hoogeboom2021argmax} with self-conditioning~\citep{chen2022analog}, enabling bidirectional long-range dependencies.
The B/Y-ion pair bias adapts AlphaFold3's pair-representation attention~\citep{abramson2024alphafold3} to fragment ion complementarity.

%% ─── 9. Conclusion ───────────────────────────────────────────────────────────
\section{Conclusion}

We presented a self-conditioned absorbing diffusion model that reaches 76.00\%~AA recall and
59.96\%~peptide accuracy on the E.\,coli EV benchmark, beating InstaNovo by +3.1~pp and +26.9~pp
(stable across 3 seeds, $\pm$0.19~pp AA).
The gains trace to two components: the PeakEncoder with B/Y-ion pair bias (+30~pp AA) and
self-conditioning (+14.6~pp peptide accuracy).
We also identified 6 structurally plausible wastewater peptides at 5\%~FDR without any reference
proteome.
Future work: larger multi-species benchmarks, BLAST confirmation of wastewater hits, and
incorporating mass constraints into training rather than inference.

%% ─── Team Contributions ──────────────────────────────────────────────────────
\section*{Team Contributions}
\textbf{Vaishak Girish Kumar:} diffusion backbone, PeakEncoder, self-conditioning, training/ablation experiments, ESM-2 validation, writing lead.
\textbf{Akshay Mohan Revankar:} CFID/SGIR reranking, beam search, mass-constraint ablations, wastewater pipeline, GitHub.
\textbf{Sanika Vilas Nanjan:} mzML preprocessing (pyOpenMS), FDR filtering, wastewater quality assessment, figures.

%% ─── LLM Usage ───────────────────────────────────────────────────────────────
\section*{LLM Usage Statement}
LLMs (Claude, GPT-4) were used only for proofreading, LaTeX formatting fixes, and library API questions.
All scientific content, model design, experimental results, and final code were produced by the authors.

%% ─── References ──────────────────────────────────────────────────────────────
\bibliographystyle{plainnat}
{\sloppy
\begin{thebibliography}{99}

\bibitem[Abramson et~al.(2024)]{abramson2024alphafold3}
Abramson, J. et~al.\ (2024).
Accurate structure prediction of biomolecular interactions with AlphaFold~3.
\textit{Nature}.

\bibitem[Austin et~al.(2021)]{austin2021structured}
Austin, J., Johnson, D.\,D., Ho, J., Tarlow, D., van den Berg, R.\ (2021).
Structured denoising diffusion models in discrete state-spaces.
\textit{NeurIPS}.

\bibitem[Chen et~al.(2022)]{chen2022analog}
Chen, T., Zhang, R., Hinton, G.\ (2022).
Analog bits: Generating discrete data using diffusion models with self-conditioning.
\textit{arXiv:2208.04202}.

\bibitem[Hoogeboom et~al.(2021)]{hoogeboom2021argmax}
Hoogeboom, E., Nielsen, D., Jaini, P., Forré, P., Welling, M.\ (2021).
Argmax flows and multinomial diffusion.
\textit{NeurIPS}.

\bibitem[Eloff et~al.(2025)]{instanovo2024}
Eloff, K. et~al.\ (2025).
InstaNovo enables diffusion-powered de novo peptide sequencing in large-scale
proteomics experiments.
\textit{Nature Machine Intelligence}, 7, 565--579.

\bibitem[Jumper et~al.(2021)]{jumper2021highly}
Jumper, J. et~al.\ (2021).
Highly accurate protein structure prediction with AlphaFold.
\textit{Nature}, 596, 583--589.

\bibitem[Lin et~al.(2023)]{lin2023evolutionary}
Lin, Z. et~al.\ (2023).
Evolutionary-scale prediction of atomic-level protein structure with a language model.
\textit{Science}, 379(6637), 1123--1130.

\bibitem[Yilmaz et~al.(2022)]{casanovo2023}
Yilmaz, M., Fondrie, W.\,E., Bittremieux, W., Oh, S., Noble, W.\,S.\ (2022).
De novo mass spectrometry peptide sequencing with a transformer model.
\textit{ICML}. (Journal version: \textit{Nature Communications}, 2024.)

\bibitem[Kantroo et~al.(2024)]{kantroo2024}
Kantroo, M., Tauber, J., Bhatt, D.\ (2024).
Pseudo-perplexity in one fell swoop for protein fitness estimation.
\textit{PRX Life}.

\bibitem[Tran et~al.(2017)]{tran2017deepnovo}
Tran, N.\,H., Zhang, A., Xin, L., Shan, B., Li, M.\ (2017).
De novo peptide sequencing by deep learning.
\textit{Proceedings of the National Academy of Sciences}, 114(31), 8247--8252.

\bibitem[Zhang et~al.(2025)]{zhang2025primenovo}
Zhang, X. et~al.\ (2025).
$\pi$-PrimeNovo: an accurate and efficient non-autoregressive deep network
for de novo peptide sequencing.
\textit{Nature Communications}, 16, 267.

\bibitem[Zhou et~al.(2024)]{zhou2024novobench}
Zhou, J. et~al.\ (2024).
NovoBench: benchmarking de novo sequencing methods in proteomics.
\textit{NeurIPS}.

\end{thebibliography}
}

\end{document}
